% Paper outline (v0.2) for collaborators. ICLR 2027 template, no title block.
% Compile: pdflatex.
\documentclass{article}
\usepackage{iclr2027_conference,times}
\iclrfinalcopy
\usepackage{amsmath}
\usepackage{enumitem}
\usepackage[table]{xcolor}

\newcommand{\TODO}{\colorbox{red!15}{\scriptsize\textbf{TODO}}}

\begin{document}

\paragraph{Working title.}
\emph{Muscle Memory for Robot Policies: Caching a VLA's Own Experience in Its
Own Representations.}

\paragraph{Outline (v0.2).}
Page budgets follow the 9-page main-text limit. The current draft policy is to
write every section in full and cut at writing time; the budget is exceeded
deliberately. Experiment codes (E0--E12) refer to the experiment list.

\begin{enumerate}[leftmargin=1.6em, itemsep=0.5em, topsep=0.3em]
\item \textbf{Introduction} ($\sim$1.25 pp; Fig.~1).
(i) Deployed robots repeat themselves while VLAs recompute every step; serving
systems schedule this demand, none of them remove it.
(ii) The central question: how much of closed-loop VLA inference is redundant,
in the inter-frame-compression sense?
(iii) Key distinction from prior retrieval caches: keys are extracted from the
teacher's own forward pass rather than from an external encoder.
(iv) System summary: training-free threshold dispatch over a library of the
teacher's own successes.
(v) Three contributions: (a) definition and first closed-loop measurement of
recoverable redundancy $R(\varepsilon)$, including its temporal structure;
(b) the serving primitive, orthogonal to distillation, quantization, and
step-reduction; (c) the frontier and its economics (E1 numbers; cross-scene
transfer \TODO). Fig.~1 combines the system diagram, the four-stage lifecycle,
and an episode hit-strip.

\item \textbf{Related work} ($\sim$0.6 pp).
Six families: robot retrieval caches (RT-Cache; key provenance is the
distinction); VLA acceleration primitives (orthogonal); the 2026
adaptive-computation line (SkiP, ElegantVLA: mechanisms without a defined
quantity or a library); policy serving (ROSA, Armory: scheduling versus
removal of demand); the cache lineage across domains (exact KV/prefix reuse,
approximate semantic reuse, closed-loop approximate reuse here); concept
ancestors (event-triggered control; uncertainty quantification and failure
detection). Closes with the first-claim: to our knowledge, the first formal definition
and closed-loop measurement of the recoverable redundancy of VLA inference;
the systems-level distinction is stated structurally (keys before the
backbone; concurrent caches key on backbone outputs and must run it every
step).

\item \textbf{Recoverable redundancy} ($\sim$0.75--1 pp).
The formalism of the companion note: step-level definition, failure of
composition, $R_{\mathcal{C}}(\varepsilon)$, the information-set
stratification, attribution, monotonicity, the intrinsic limit. Full proofs
and extended remarks in Appendix~A. Three interfaces to the rest of the paper
are stated explicitly: calibration solves Eq.~(3) empirically; the E1 sweep is
the empirical Pareto frontier of the threshold class; the attribution
quadruple organizes the experimental chapters.

\item \textbf{System} ($\sim$1.5 pp).
Key extraction (tap point, layout, pooling; the corresponding tap point in a
second architecture); library construction under a success filter; retrieval
with per-field normalization and weighted fusion; threshold calibration as the
empirical program of Eq.~(3); the three execution tiers (replay; warm start as
experience-guided step reduction; full inference); the deployment lifecycle
(collect, shadow, calibrate, run and append); the cost account (E0: hit
$\approx$ 14\,ms against 72\,ms, the $\mathrm{IR}$ floor).

\item \textbf{Measuring $R(\varepsilon)$} ($\sim$1.5 pp; Fig.~2, Fig.~3,
Table~1).
(5.1) Main frontier: four panels, two benchmarks $\times$ two teachers, each
panel with the threshold sweep (E1) and the oracle arm (E2 \TODO) overlaid;
teacher anchors marked.
(5.2) Library-size scaling (E3).
(5.3) Temporal structure (E4), one half-page: stickiness, run lengths,
lockstep; one sentence on score discriminability (E5 \TODO).
(5.4) Key-source ablation (E7a \TODO): internal representations against
external encoders, the direct evidence for the title claim; extraction depth
and fusion (E7b, E7c) remain in Appendix~E.
Table~1: pre-registered operating points with success rate, replay share,
GPU-time per step, robots-per-GPU ($1/\mathrm{IR}$), and resident memory;
anchored by the serving mini-bench (E11 \TODO).

\item \textbf{Does experience transfer?} ($\sim$0.75--1 pp; Fig.~4).
Cross-scene transfer on RoboCasa365 (E6 \TODO): library built in one kitchen,
queried in held-out kitchens, two teachers; analysis stratified by task type
(preliminary: large-motion tasks transfer, precise dynamic tasks do not).
Fallback if the data are not ready: the main figure retreats to LIBERO and
this section absorbs all RoboCasa material.

\item \textbf{Real-robot validation} ($\sim$0.3--0.5 pp; E13 \TODO).
The full deployment lifecycle executed on a physical SO-101 arm: collect,
shadow, calibrate, serve. Three tasks stratified by motion scale (drawer or
lid, pick-and-place, insertion), approximately 330 episodes, paired evaluation
at two operating points. Reports success rate, replay share, and end-to-end
latency, where a hit is a local replay and a miss is a network round trip plus
inference; the latency column instantiates the per-robot axis with physical
measurements. Positioned as an existence proof, not a main result.

\item \textbf{Alternatives} ($\sim$0.75 pp; Table~2).
(7.1) Distilled students: the regime map; the acquisition-cost column; the
success-rate loss reported without qualification (E9).
(7.2) Trained routers: counterfactual labels, retraining coupling to the
library, the threshold as a byproduct of retrieval; the comparison E10.
(7.3) Cheaper teachers: orthogonality and composability; warm start against
blind step reduction (E8); appendix for the three-way control.

\item \textbf{Limitations} ($\sim$0.3 pp).
Real-robot validation is small-scale and serves as an existence proof;
quasi-static tasks (staleness sensitivity of replay); single embodiment
class; the robots-per-GPU projection is an idealization anchored at one
measured point; library growth is append-only with no eviction policy.

\item \textbf{Conclusion} ($\sim$0.1 pp).
\end{enumerate}

\paragraph{Floats.}
Fig.~1 system, lifecycle, hit-strip. Fig.~2 main frontier, four panels with
oracle overlay. Fig.~3 library-size scaling and temporal structure. Fig.~4
cross-scene transfer. Fig.~5 real-robot latency and success (E13). Table~1 operating points and economics. Table~2
alternatives. One comparison table in Section~2.

\paragraph{Appendix plan.}
A: full formalism, proofs, extended remarks. B: cost accounting, serving
mini-bench (E11), warm-start three-way control (E8). C: history-augmented
negatives, one table (E12). D: calibration protocol, pre-registration records,
statistics. E: per-task breakdowns, key-space ablations E7b (depth) and E7c (fusion),
score discriminability (E5). F: implementation, reproducibility, data and code
release. A further appendix section, \emph{Design choices and alternatives
considered}, answers anticipated questions in declarative form (learned
routers, distillation, memorization of the evaluation set, unbounded library
growth, safety of replayed actions), restricted to questions with substantiated
answers.

\paragraph{Remaining collection.}
E2 (oracle arm), E11 (serving mini-bench), E8 refresh, E13 (SO-101; teacher
fine-tuning is the prerequisite); E1 RoboCasa cells, E6, and E10 in progress;
E5 and E7 offline.

\end{document}
